% Compile twice with pdflatex.  This source is self-contained; no BibTeX run
% or external figure files are needed.  See README.md for draft status.
\documentclass[11pt]{article}
\usepackage[T1]{fontenc}
\usepackage[utf8]{inputenc}
\usepackage{lmodern}
\usepackage[letterpaper,margin=1in]{geometry}
\usepackage{amsmath,amssymb,amsthm,mathtools}
\usepackage{microtype}
\usepackage{enumitem}
\usepackage{booktabs}
\usepackage{etoolbox}
\usepackage[hidelinks,pdfencoding=auto]{hyperref}
\usepackage[capitalise,nameinlink,noabbrev]{cleveref}
\hypersetup{
  pdftitle={Square-Root Circuit Bounds for Iterated Functions},
  pdfauthor={Samuel Schlesinger},
  pdfsubject={Preliminary research draft: nonuniform Boolean circuit upper bounds},
  pdfkeywords={circuit complexity, function iteration, functional graphs, nonuniform computation, sorting networks, permutation inversion, one-way functions}
}
\newtheorem{theorem}{Theorem}[section]
\newtheorem{lemma}[theorem]{Lemma}
\newtheorem{proposition}[theorem]{Proposition}
\newtheorem{corollary}[theorem]{Corollary}
\theoremstyle{definition}
\newtheorem{definition}[theorem]{Definition}
\theoremstyle{remark}
\newtheorem{remark}[theorem]{Remark}
% Keep theorem statements together while allowing their proofs to break.
\newcommand{\keepstatement}[1]{%
  \BeforeBeginEnvironment{#1}{\begin{samepage}}%
  \AfterEndEnvironment{#1}{\end{samepage}}}
\keepstatement{theorem}
\keepstatement{lemma}
\keepstatement{proposition}
\keepstatement{corollary}
\keepstatement{definition}
\newcommand{\bits}{\{0,1\}}
\newcommand{\CC}{\mathrm{C}}
\newcommand{\Iter}{\operatorname{Iter}}
\newcommand{\lcm}{\operatorname{lcm}}
\newcommand{\poly}{\operatorname{poly}}
\newcommand{\ord}{\operatorname{ord}}
\newcommand{\floor}[1]{\left\lfloor #1\right\rfloor}
\newcommand{\ceil}[1]{\left\lceil #1\right\rceil}
\DeclareMathOperator*{\argmin}{arg\,min}
\setlist[enumerate]{itemsep=3pt,topsep=5pt}
\setlength{\emergencystretch}{2em}
\numberwithin{equation}{section}
\title{Square-Root Circuit Bounds\\for Iterated Functions}
% Author field is provisional and editable; no affiliation is assumed.
\author{Samuel Schlesinger}
\date{September 29, 2026\\[3pt]\small Preliminary draft}

\begin{document}
\maketitle

\begin{abstract}
Let $F:\bits^n\to\bits^n$ be computed by a Boolean circuit of size $s$, and
write $A=s+n^3$.  We prove that every fixed iterate $F^T$, for an arbitrary
integer $T\geq 0$, has a bounded-fan-in Boolean circuit of size
\[
  O\!\left(A+\sqrt{n^3 2^n A}\right).
\]
The circuit is nonuniform in $F$ and $T$, but the bound is independent of the
magnitude and bit length of $T$.  No injectivity assumption is required.
Consequently, all fixed iterates of a polynomial-size update rule have
circuits of size $2^{n/2+O(\log n)}$.  The same exponential bound holds when
an $n$-bit iteration count is part of the input.  We also give a batched
construction that shares preprocessing across independent trajectories.

The proof combines sparse checkpoints in a functional graph with a
forward-only correction identity and batched dictionary lookup implemented
by sorting networks.  All lookup and selection gates are included in the
size bound.  The theorem recovers square-root-size circuits for permutation
inversion, consistently with possible one-wayness against polynomial-size
adversaries.  Finally, for even $n$ we construct an explicit linear-size
permutation on $n$ bits for which some fixed power requires
$\Omega(2^{n/2}/n)$ gates.  Such hard powers necessarily have binary
exponents of length $\Omega(2^{n/2}/n^2)$.  The counting argument neither
identifies an explicit hard power nor establishes one-way functions.
\end{abstract}

\noindent\textbf{Keywords.} Boolean circuit complexity; function iteration;
functional graphs; nonuniform computation; sorting networks; permutation
inversion; one-way functions.

\section{Introduction}\label{sec:intro}

A simple transition rule can describe a computation of enormous duration.
How complicated can its final-state map be as a Boolean function of the
initial state?  We study this question for a fixed map
$F:\bits^n\to\bits^n$.  The computational resource is the total number of
Boolean gates, rather than the number of evaluations of $F$, the number of
memory probes, or the time of a random-access algorithm.

Let $N=2^n$ and suppose that $F$ has a circuit of size $s$.  Directly
unrolling $T$ applications gives a circuit of size at most $Ts$ for $F^T$.
Alternatively, hardwiring a complete truth table and testing the input
against all $N$ states gives an $O(nN)$-gate circuit, independently of $T$.
The latter bound ignores the possibility that $F$ itself is much simpler
than an arbitrary map on its state space.  Our main result exploits that
possibility.

\begin{theorem}[Fixed-power tradeoff]\label{thm:main}
Let $n\geq 1$, let $F:\bits^n\to\bits^n$ satisfy $\CC(F)\leq s$, and put
$N=2^n$ and $A=s+n^3$.  For every integer $T\geq 0$ and every integer
$1\leq L\leq N$,
\begin{equation}\label{eq:tradeoff}
  \CC(F^T)
  =O\!\left(Ls+n^3\left(L+\frac{N}{L}\right)\right).
\end{equation}
In particular,
\begin{equation}\label{eq:optimized}
  \CC(F^T)
  =O\!\left(A+\sqrt{n^3NA}\right).
\end{equation}
The implicit constants are absolute.  The circuit may depend on $F$ and
$T$, but the size bound does not depend on the bit length of $T$.
\end{theorem}

Here and throughout, circuits use the full fan-in-two Boolean basis and
have unrestricted fan-out.  The results also hold over the basis
$\{\mathsf{AND},\mathsf{OR},\mathsf{NOT}\}$ with a constant-factor change.
Circuit size refers to the whole multi-output circuit, not to a separate
size bound for each output bit.

\begin{corollary}\label[corollary]{cor:poly}
If $\CC(F_n)=n^{O(1)}$, then
\begin{equation}\label{eq:half}
  \sup_{T\geq 0}\CC(F_n^T)\leq 2^{n/2+O(\log n)}.
\end{equation}
The constant hidden in $O(\log n)$ may depend on the polynomial bound for
$\CC(F_n)$, but not on $T$.
\end{corollary}

The quantifiers in \cref{thm:main} are important.  For each $F$ and $T$ we
construct a circuit that works for \emph{every} initial state.  We do not
construct a single circuit taking an arbitrary-length exponent as input.
Nevertheless, an $m$-bit exponent can be supplied as input at a multiplicative
cost of $m$:
\begin{equation}\label{eq:intro-variable}
  \CC\bigl((x,t)\mapsto F^t(x)\bigr)
  =O\!\left(m\left(A+\sqrt{n^3NA}\right)\right),
  \qquad t\in\bits^m.
\end{equation}
We prove this in \cref{sec:variable}.  In particular, $m=n$ preserves the
exponent in \cref{eq:half}.  We also give a shared-dictionary construction
for $r$ independent initial states, replacing the main square-root term
$r\sqrt{n^3NA}$ of separate constructions by $\sqrt{rn^3NA}$; see
\cref{sec:batch}.

For arbitrary fixed powers, the exponent $1/2$ cannot be improved in
general.  In \cref{sec:lower} we give an explicit family of permutations
$P_n$ with $\CC(P_n)=O(n)$ and
\begin{equation}\label{eq:intro-lower}
  \max_{T\geq 0}\CC(P_n^T)=2^{n/2+O(\log n)}
\end{equation}
for even $n$.  The hard exponent is obtained by counting over a very large
cyclic group of powers.  This is not an explicit exponential circuit lower
bound: the permutation is explicit, but the hard power is not specified
by an efficient procedure.  Moreover, the same permutation has an easy
iterator when the time parameter has polynomially many bits.

For a permutation $P$, choosing $T=\ord(P)-1$ recovers a circuit for
$P^{-1}$.  Thus every polynomial-size permutation on $n$ bits admits an
exact inverter of size $2^{n/2+O(\log n)}$.  This exponential upper bound
is compatible with one-wayness against polynomial-size circuits.
We explain this consequence and its limitations in \cref{sec:one-way}.

\subsection{Proof overview}

The functional graph of $F$ consists of directed cycles with directed
trees feeding into them.  For a parameter $L$, we select $O(N/L)$
checkpoints so that every trajectory either encounters one within $2L$
steps or repeats a state within that many steps.  A repeated state reveals
a period and permits direct recovery of the answer from the initial
trajectory.

Set $K=2L$ and assume $T\geq K$.  At each checkpoint $h$, store
$z_h=F^{T-K}(h)$.  If the initial trajectory encounters $h$ at time $d\leq K$,
then
\begin{equation}\label{eq:overview-identity}
  F^{K-d}(z_h)
  =F^{K-d}\bigl(F^{T-K}(F^d(x))\bigr)
  =F^T(x).
\end{equation}
Thus only a short \emph{forward} correction is needed after the lookup.
There is no inverse operation, so the argument applies to noninjective
maps.

A separate lookup for each of $\Theta(L)$ visited states against
$\Theta(N/L)$ checkpoints would cost $\Theta(N)$ comparisons.  Instead,
we generate the full initial trajectory before looking anything up and
resolve all its queries together.  Sorting the queries and dictionary
records reduces this cost to
$\poly(n)(L+N/L)$ Boolean gates.  Combining the initial and final short
simulations with this batched lookup gives \cref{eq:tradeoff}.

\subsection{Related work and status of the formulation}\label{sec:related}

Preprocessing trajectories to exchange storage for later computation is
classical.  Hellman's cryptanalytic time--memory tradeoff is an early
example of this approach~\cite{Hellman80}.  More directly, Munro, Raman,
Raman, and Rao~\cite[Section~3.1]{MRRR12} give a shortcut representation
for permutation inversion: with a parameter $L$, their black-box
construction uses $O(L)$ evaluations and
$O((N/L)(\log N+\log L))$ auxiliary bits.  They also give succinct
representations of general functions supporting positive powers with a
word-size exponent~\cite[Section~5]{MRRR12}.  Those results concern
representation space, RAM time, and black-box access, rather than total
Boolean gate count.

A closer circuit antecedent is Hirahara, Ilango, and
Williams~\cite[Section~3.2]{HIW24}: they obtain square-root-size
permutation-inversion circuits using checkpoints and batched lookup.
Their lookup theorem builds on a lemma of Paul.  This draft generalizes
that permutation-inversion construction to arbitrary fixed forward powers
of possibly noninjective maps.  It does not subsume their general-function
inversion and compression theorems.

Eppstein~\cite{Eppstein23} studies uniform iteration of polynomial-time
bijections, with the iteration count supplied in binary, and characterizes
the resulting class of functional problems as
$\mathsf{FP}^{\mathsf{PSPACE}}$.  Our resource question differs: the map
is fixed, the circuits are nonuniform, and the size bound remains
exponential in the state length.  Our result supplies no polynomial-time
algorithm for that iteration problem.

The present construction uses the same checkpoint and batched-lookup
principles.  Batcher's sorting networks~\cite{Batcher68} implement the
lookup, while the shifted payload in \cref{eq:overview-identity} and the
checkpoint-or-repetition lemma handle general powers and functional graphs.
No RAM memory probe is treated as a unit-cost circuit operation.
Checkpointing, sorting-based lookup, and square-root-size permutation
inversion are not claimed as new.  Historical priority for the full
fixed-power theorem and the shared bounds in this draft remains unresolved.

\section{Model and circuit primitives}\label{sec:model}

\subsection{Conventions}

All logarithms are to base two.  The notation $F^t$ denotes $t$-fold
composition, with $F^0$ equal to the identity.  The quantity $\CC(G)$ is the
minimum number of internal gates in an acyclic Boolean circuit computing
all output bits of $G$.  Each internal gate computes a Boolean function
of at most two bits.  Inputs, wires, fan-out, and output designations are
not charged as gates.  The two constant bits are available; charging a
constant number of gates for them would not change any bound.

The circuit for an iterate is allowed to depend on the entire function
$F$ and on its fixed exponent.  Preprocessing time and the time required
to find the circuit are not bounded.  Hardwired data do not provide free
addressing: every comparison, multiplexer, and routing operation used to
access them is realized with Boolean gates and included in our estimates.

All state words have $n$ bits.  Because $L\leq 2^n$, indices up to $2L$ fit
in $n+2$ bits.  In the batched construction $r\leq 2^n$, so query identifiers
and record positions also have $O(n)$ bits.  We retain explicit polynomial
factors in $n$ for a concrete implementation; these factors are not
claimed to be optimal.

\subsection{Arithmetic and selection}

Comparison, addition, subtraction, and conditional selection of $w$-bit
words have $O(w)$-gate circuits.  Binary long division gives an
$O(w^2)$-gate circuit for the remainder of a $w$-bit integer modulo a
positive $w$-bit integer: process the dividend bits from most to least
significant, and at each stage conditionally subtract the modulus from
the current shifted remainder.  There are $O(w)$ stages and $O(w)$ gates
per stage.

A circuit can select one of $q$ words of width $w$ using $O(qw)$ gates by
a binary tree of word multiplexers, padding to a power of two when
necessary.  It can also select one flagged record from an ordered list of
$q$ records in $O(qw)$ gates.  For example, a left-to-right scan maintains
a bit indicating whether a flagged record has already been found and
retains the first one.  This scan is simply an acyclic chain of gates;
no depth bound is needed here.

\begin{lemma}[Batched fixed-dictionary lookup]\label[lemma]{lem:lookup}
Suppose a fixed dictionary consists of $M$ distinct $w$-bit keys, each
with a $w$-bit payload.  Given $Q\geq 1$ query keys, a Boolean circuit can
output a match bit and the corresponding payload for every query, in the
original query order, using
\begin{equation}\label{eq:lookup-general}
  O\!\left((M+Q)\bigl(w+\log(M+Q+2)\bigr)
                    \log^2(M+Q+2)\right)
\end{equation}
gates.  The payload of an unmatched query can be set to zero.  In
particular, if $w=O(n)$ and $\log(M+Q+2)=O(n)$, the cost is
$O(n^3(M+Q))$.
\end{lemma}

\begin{proof}
Let $R=M+Q$.  Form one record for each dictionary entry and one for each
query.  A record contains a key, a dictionary/query tag, a payload field,
and an original-position field.  The record width is
$O(w+\log(R+2))$.

Sort the records lexicographically by key, then by the tag with dictionary
records first, and then by original position.  A sorting network with
$O(R\log^2(R+2))$ comparators exists~\cite{Batcher68}.  A comparator
and the resulting conditional swap of the entire records cost
$O(w+\log(R+2))$ Boolean gates.  Padding uses a separate dummy-record bit
so that no actual key is reserved as a sentinel.

In the sorted list, a dictionary record precedes every query with its key.
Scan the list, resetting the match bit when the key changes and propagating
the payload from a dictionary record through the equal-key run that
follows it.  If an equal-key run has no dictionary record, every query in
that run receives match bit zero.  Dictionary keys are distinct, so no
run contains two competing payloads.  This scan uses
$O(R(w+\log(R+2)))$ gates.

Finally, sort records by a tag that places queries first and then by the
query's original position.  The first $Q$ records give the requested
outputs in order.  A second sorting network has the same asymptotic cost
as the first.  Every operation described is a fixed acyclic circuit.
\end{proof}

\begin{lemma}[Detecting a repeated state]\label[lemma]{lem:repeat}
Given $Q$ words of $w$ bits, a circuit can return a flag indicating whether
two words are equal and, when the flag is one, indices $i<j$ of two equal
words.  Its size is
\[
  O\!\left(Q\bigl(w+\log(Q+2)\bigr)\log^2(Q+2)\right).
\]
In particular, for $Q\leq 2^{n+O(1)}$ and $w=O(n)$, the size is $O(n^3Q)$.
\end{lemma}

\begin{proof}
Attach each word's original index and sort by word and then by index.
A duplicate exists exactly when two adjacent sorted records have the
same word.  Compare adjacent words and select one flagged pair.  Sorting
by index within an equal-word run guarantees that the returned indices
satisfy $i<j$.  The remaining comparisons and selection use at most
$O(Q(w+\log(Q+2)))$ gates.
\end{proof}

\section{Sparse checkpoints in functional graphs}\label{sec:checkpoints}

The functional graph of a map $F$ has one vertex for each state and the
edge $v\to F(v)$ at every vertex.  Each weakly connected component has a
unique directed cycle, with directed trees feeding into it.  For a
noncycle vertex, its \emph{depth} is the number of steps needed to reach
the cycle.

\begin{lemma}[Checkpoint or repetition]\label[lemma]{lem:checkpoints}
Let $F$ be a function on a set of $N$ states, and let $1\leq L\leq N$.
There is a set $H$ of at most $2N/L$ states such that, for every $x$, the
list
\[
  x,F(x),\ldots,F^{2L}(x)
\]
either contains an element of $H$ or contains a repeated state.
\end{lemma}

\begin{proof}
Let $N_{\mathrm{tree}}$ and $N_{\mathrm{cycle}}$ be the numbers of noncycle
and cycle vertices.  Among the $L$ residue classes of noncycle depths
modulo $L$, choose one of minimum cardinality, and put every vertex in
that class into $H$.  This adds at most $N_{\mathrm{tree}}/L$ vertices.
The choice is global across all components; choosing a residue separately
for every tree is unnecessary.

For each cycle of length $\ell\geq L$, number its vertices consecutively
in the forward direction and mark vertices with indices
$0,L,2L,\ldots$ below $\ell$.  Consecutive marked vertices, including the
last and first around the cycle, are at forward distance at most $L$.
The number of marks on that cycle is
\[
  \ceil{\ell/L}\leq 2\ell/L.
\]
Leave cycles of length less than $L$ unmarked.  The total number of marks
is at most
\[
  \frac{N_{\mathrm{tree}}}{L}
    +\frac{2N_{\mathrm{cycle}}}{L}
  \leq \frac{2N}{L}.
\]

If the starting depth is at least $L$, the first $L$ depths encountered
include every residue class modulo $L$, so a marked noncycle vertex is
encountered.  If the starting depth is less than $L$, the trajectory
reaches a cycle in fewer than $L$ steps.  On a cycle of length at least
$L$, it reaches a marked vertex in at most another $L$ steps.  On a
shorter cycle, it repeats a state in fewer than another $L$ steps.  The
same cycle analysis applies to a trajectory starting on a cycle.  All
cases are covered by the displayed list.
\end{proof}

\begin{remark}
It is essential that the lemma allows a short repetition instead of a
checkpoint.  A function can have $N$ fixed points, in which case a set
meeting every eventual cycle must contain all $N$ states.  Such cycles
need no stored shortcut: their behavior is already visible locally.
\end{remark}

\section{Circuits for fixed powers}\label{sec:fixed}

We now prove \cref{thm:main}.  The construction is uniform as a description
of a circuit from the chosen constants, but no efficient procedure for
computing those constants is asserted.

\subsection{Short trajectories and shifted payloads}

Fix $L$ and let $K=2L$.  If $T<K$, use $T$ copies of a circuit for $F$;
this costs at most $2Ls$ gates.  Hence assume $T\geq K$.
Choose $H$ as in \cref{lem:checkpoints}.  For every $h\in H$, hardwire the
$n$-bit payload
\begin{equation}\label{eq:payload}
  z_h=F^{T-K}(h).
\end{equation}
The dictionary is fixed independently of the input $x$ and has at most
$2N/L$ entries.

On input $x$, first compute
\begin{equation}\label{eq:prefix}
  x_0=x,\quad x_1=F(x),\quad\ldots,\quad x_K=F^K(x).
\end{equation}
This requires $Ks$ gates.  Apply \cref{lem:lookup} to all $K+1$ states in
\cref{eq:prefix}, using the dictionary $h\mapsto z_h$.  Select one
matching record, if any, and retain its original trajectory index $d$.
This gives a match flag and a pair $(d,z_h)$ with $x_d=h$ whenever the
flag is one.

For the matched case, compute a second short trajectory
\[
  z_h,F(z_h),\ldots,F^K(z_h)
\]
and select its entry at position $K-d$.  Its value is
\begin{align}
  F^{K-d}(z_h)
    &=F^{K-d}\bigl(F^{T-K}(h)\bigr)\notag\\
    &=F^{T-d}(h)\notag\\
    &=F^{T-d}\bigl(F^d(x)\bigr)
     =F^T(x).\label{eq:identity}
\end{align}
All exponents used here are nonnegative.  This second trajectory costs
another $Ks$ gates; selection costs $O(nK)$ gates.

\subsection{The repeated-state case and very large exponents}

Use \cref{lem:repeat} on \cref{eq:prefix}.  If no checkpoint was found,
\cref{lem:checkpoints} guarantees indices
\[
  0\leq i<j\leq K,\qquad x_i=x_j.
\]
Put $p=j-i$.  Determinism implies $x_{u+p}=x_u$ for all $u\geq i$.
The number $p$ need not be the least period.  Since $T\geq K\geq i$,
\begin{equation}\label{eq:period}
  F^T(x)=x_{\,i+((T-i)\bmod p)}.
\end{equation}
The selected index is at most $i+p-1=j-1$, so the required state already
occurs in the initial trajectory.

To implement \cref{eq:period} without paying for the bit length of $T$,
hardwire the table
\begin{equation}\label{eq:residue-table}
  R_T(p)=T\bmod p,\qquad 1\leq p\leq K.
\end{equation}
There are $K$ entries, each of $O(n)$ bits.  A multiplexer tree retrieves
$R_T(p)$ in $O(nK)$ gates.  Compute $i\bmod p$ by binary long division
and then compute
\[
  (T-i)\bmod p
    =\bigl(R_T(p)-(i\bmod p)\bigr)\bmod p
\]
using a comparison, subtraction, and a conditional addition of $p$.
These operations take $O(n^2)$ gates.  A further $O(nK)$ gates select the
state prescribed in \cref{eq:period}.

A circuit evaluates both branches, so all intermediate operations must
also be defined when their output will be discarded.  If a repeated pair
is absent, supply the dummy values $i=0$ and $p=1$ to the arithmetic
branch.  If a checkpoint is absent, supply a dummy state and $d=0$ to
the second trajectory.  The final output selects the checkpoint branch
when its match flag is one, and otherwise selects the repetition branch.
The checkpoint lemma ensures that the selected branch is always valid.
In particular, no division by zero occurs on any input.

\subsection{Size bound}

The two simulations together use at most $2Ks=4Ls$ gates.  The dictionary
and query list contain
\[
  |H|+(K+1)=O\!\left(\frac{N}{L}+L\right)
\]
records of $O(n)$ bits, and the logarithm of this number is $O(n)$.
Thus \cref{lem:lookup} costs
$O(n^3(N/L+L))$ gates.  Repeated-state detection costs $O(n^3L)$ gates,
and all remaining lookup, arithmetic, and selection operations cost
$O(nL+n^2)$ gates.  Since $n,L\geq 1$, the latter cost is absorbed by
$O(n^3L)$.  This proves \cref{eq:tradeoff}.

Write the tradeoff as $O(AL+n^3N/L)$, where $A=s+n^3$, and choose
\begin{equation}\label{eq:L-choice}
  L=\max\left\{1,\ceil{\sqrt{n^3N/A}}\right\}.
\end{equation}
Since $A\geq n^3$, this choice satisfies $L\leq\ceil{\sqrt N}\leq N$.
If $A\leq n^3N$, then
$\sqrt{n^3N/A}\leq L\leq 2\sqrt{n^3N/A}$.  Hence
\[
  AL=O(\sqrt{n^3NA}),\qquad
  n^3N/L=O(\sqrt{n^3NA}).
\]
If $A>n^3N$, the choice is $L=1$ and the cost is $O(A)$.
This establishes
\cref{eq:optimized} and completes the proof of \cref{thm:main}.
\hfill$\square$

\begin{remark}[Where the exponent is stored]\label[remark]{rem:nonuniform}
For $T\geq K$, the information about $T$ used by the circuit is confined
to the checkpoint payloads in \cref{eq:payload} and the residue table in
\cref{eq:residue-table}.  Their total length is
$O(n(N/L+L))$ bits.  The construction does not place the binary expansion
of $T$ on a long hidden input tape, and it does not perform uncharged
arithmetic on an arbitrarily large integer.  Computing the hardwired
information is a preprocessing task outside the asserted circuit bound.
\end{remark}

\section{Input times and shared evaluation}

\subsection{An iteration count supplied as input}\label{sec:variable}

\begin{definition}
For $m\geq 1$, define
\[
  \Iter_{F,m}:\bits^n\times\bits^m\longrightarrow\bits^n,
  \qquad \Iter_{F,m}(x,t)=F^{\operatorname{val}(t)}(x),
\]
where $\operatorname{val}(t)$ is the nonnegative integer encoded by $t$.
\end{definition}

\begin{theorem}[Variable-time iteration]\label{thm:variable}
Under the assumptions of \cref{thm:main}, for every $m\geq 1$,
\begin{equation}\label{eq:variable}
  \CC(\Iter_{F,m})
  =O\!\left(m\left(A+\sqrt{n^3NA}\right)\right).
\end{equation}
In particular,
\[
  \CC(\Iter_{F,n})
  =O\!\left(n\left(A+\sqrt{n^3NA}\right)\right).
\]
If $s$ and $m$ are polynomial in $n$, then
$\CC(\Iter_{F,m})\leq 2^{n/2+O(\log n)}$.
\end{theorem}

\begin{proof}
Write $t=\sum_{j=0}^{m-1}t_j2^j$.  For each $j$, apply
\cref{thm:main} directly to the original map $F$ and the fixed exponent
$2^j$, obtaining a circuit for $G_j=F^{2^j}$.  Starting with $v_0=x$, use
\[
  v_{j+1}=
  \begin{cases}
    G_j(v_j),&t_j=1,\\
    v_j,&t_j=0.
  \end{cases}
\]
Inductively,
$v_j=F^{\sum_{a<j}t_a2^a}(x)$, so $v_m=F^t(x)$.
There are $m$ circuits of the size in \cref{eq:optimized} and $m$
$n$-bit multiplexers.  The extra $O(mn)$ gates are absorbed by the claimed
bound.  Each $G_j$ is constructed from $F$ using the theorem; it is not
obtained by duplicating the circuit for $G_{j-1}$.
\end{proof}

\subsection{Sharing a dictionary among independent trajectories}
\label{sec:batch}

Let
\[
  B_{F,T,r}(x_1,\ldots,x_r)
   =\bigl(F^T(x_1),\ldots,F^T(x_r)\bigr).
\]
The inputs $x_1,\ldots,x_r$ are independent $n$-bit words.  The following
construction uses one checkpoint dictionary for all of them.

\begin{theorem}[Batched fixed powers]\label{thm:batch}
Under the assumptions of \cref{thm:main}, for every $1\leq r\leq N$,
$T\geq 0$, and $1\leq L\leq N$,
\begin{equation}\label{eq:batch-tradeoff}
  \CC(B_{F,T,r})
   =O\!\left(rLs+n^3\left(rL+\frac{N}{L}\right)\right).
\end{equation}
Consequently,
\begin{equation}\label{eq:batch-opt}
  \CC(B_{F,T,r})
   =O\!\left(rA+\sqrt{rn^3NA}\right).
\end{equation}
\end{theorem}

\begin{proof}
Set $K=2L$.  If $T<K$, unroll $T$ copies of $F$ independently for each
input.  Otherwise use the checkpoint set and dictionary from
\cref{sec:fixed}.  Generate the first $K+1$ states of all $r$ trajectories.
Together with one copy of the dictionary, these give
\[
  R=|H|+r(K+1)=O(N/L+rL)
\]
records.  Attach the trajectory number and within-trajectory position to
each query.  Because $r\leq N$ and $L\leq N$, the record width and
$\log(R+2)$ are $O(n)$.  One application of \cref{lem:lookup} therefore
resolves all queries in $O(n^3(N/L+rL))$ gates and returns them in their
original order.

Select a matching checkpoint separately within each trajectory.  Run
repeated-state detection separately on each trajectory as well; this
costs $O(rn^3L)$ gates and does not confuse a collision between two
different trajectories with a repetition along one trajectory.
The correction trajectory and output selection are also performed
separately for each input.  The fixed residue table may be implemented
with a separate $O(nL)$-gate lookup for each trajectory, whose total
cost is absorbed by $O(rn^3L)$.  The two simulation phases together cost
$O(rLs)$ gates.  This proves \cref{eq:batch-tradeoff}.

\begin{samepage}
Choose
\[
  L=\max\left\{1,\ceil{\sqrt{n^3N/(rA)}}\right\}.
\]
\end{samepage}
As before $1\leq L\leq N$.  If $rA\leq n^3N$, then
$rAL=O(\sqrt{rn^3NA})$ and $n^3N/L=O(\sqrt{rn^3NA})$.
If $rA>n^3N$, the choice is $L=1$ and the cost is $O(rA)$.
This gives \cref{eq:batch-opt}.
\end{proof}

\begin{corollary}[Independent input times]\label[corollary]{cor:batch-variable}
For $1\leq r\leq N$ and $m\geq 1$, the map
\[
  \bigl((x_1,t_1),\ldots,(x_r,t_r)\bigr)
   \longmapsto
  \bigl(F^{t_1}(x_1),\ldots,F^{t_r}(x_r)\bigr),
  \qquad t_i\in\bits^m,
\]
has a circuit of size
\begin{equation}\label{eq:batch-variable}
  O\!\left(m\left(rA+\sqrt{rn^3NA}\right)\right).
\end{equation}
\end{corollary}

\begin{proof}
Use the binary-power construction of \cref{thm:variable}.  At stage $j$,
apply \cref{thm:batch} with update map $F$ and exponent $2^j$ on the $r$
current states, then use the $j$th bit of each $t_i$ to select whether to
retain or update that state.  The $O(mrn)$ selection gates are absorbed
by \cref{eq:batch-variable}.
\end{proof}

For polynomial $s$ and $m$, and $r=\floor{2^{\alpha n}}$ with fixed
$0\leq\alpha\leq 1$, \cref{eq:batch-variable} is at most
\[
  2^{(1+\alpha)n/2+O(\log n)}.
\]
This is an upper-bound sharing statement.  It is not a direct-sum
separation against optimal circuits for a single trajectory.

\section{A matching exponent for arbitrary fixed powers}\label{sec:lower}

This section gives an explicit linear-size permutation whose collection
of powers is large enough that some member has circuit complexity at
least $2^{n/2}/\poly(n)$.  The counting argument is elementary and does
not require an asymptotic estimate for the maximal order of a permutation.

\subsection{A permutation with cycles of all lengths up to \texorpdfstring{$2^{n/2}-1$}{2 to the power n/2, minus 1}}

Let $n=2k$ and $B=2^k-1$.  Write a state as a pair of $k$-bit integers
$(a,b)$.  Define
\begin{equation}\label{eq:permutation}
 P_n(a,b)=
 \begin{cases}
   (a,b+1\bmod a),&1\leq a\leq B\ \text{and}\ 0\leq b<a,\\
   (a,b),&\text{otherwise}.
 \end{cases}
\end{equation}
For each $1\leq a\leq B$, the states
$(a,0),(a,1),\ldots,(a,a-1)$ form one cycle of length $a$; every other
state is fixed.  Hence $P_n$ is a permutation and
\begin{equation}\label{eq:order}
  \ord(P_n)=D_B:=\lcm(1,2,\ldots,B).
\end{equation}

Despite the notation in \cref{eq:permutation}, the update does not require
a general division circuit.  Compare $b$ with $a$, increment $b$, and,
when $b<a$, replace $b$ by zero exactly when the increment equals $a$.
Otherwise retain the appropriate increment or original value.
Comparison, increment, and selection each have $O(k)$-gate circuits.
The first coordinate is passed through unchanged.  Thus
\begin{equation}\label{eq:linear-update}
  \CC(P_n)=O(n),
\end{equation}
and the family has a direct polynomial-time circuit construction.

\begin{lemma}\label[lemma]{lem:lcm}
For every integer $B\geq 1$,
\begin{equation}\label{eq:lcm-lower}
  \lcm(1,2,\ldots,B)\geq \frac{2^B}{B+1}.
\end{equation}
\end{lemma}

\begin{proof}
We show that every binomial coefficient $\binom Bj$ divides $D_B$.
For a prime $p$, factorial valuations give
\[
 \nu_p\!\binom Bj
 =\sum_{a\geq 1}
   \left(\floor{\frac{B}{p^a}}-
         \floor{\frac{j}{p^a}}-
         \floor{\frac{B-j}{p^a}}\right).
\]
Each summand is zero or one, and it is zero when $p^a>B$.
Therefore
\[
  \nu_p\!\binom Bj\leq \floor{\log_p B}=\nu_p(D_B),
\]
where both sides are zero for $p>B$.  This proves divisibility.
Since $\sum_{j=0}^B\binom Bj=2^B$, some binomial coefficient is at least
$2^B/(B+1)$, and $D_B$ is at least that large.
\end{proof}

\begin{lemma}[Counting multi-output circuits]\label[lemma]{lem:count}
The number of functions $\bits^n\to\bits^n$ computable with at most $q$
fan-in-two Boolean gates is at most
\begin{equation}\label{eq:count}
  (q+1)16^q(n+q+2)^{2q+n}
  =2^{O((q+n)\log(n+q+2))}.
\end{equation}
\end{lemma}

\begin{proof}
Topologically order the gates.  At each gate choose one of at most $16$
Boolean operations and two predecessors among the inputs, the two
constants, and gates.  Allowing all $n+q+2$ possible predecessors at each
position only increases the count.  Choose the $n$ output wires from
these same possibilities.  Sum over gate counts from zero through $q$;
each summand is bounded by the expression without its factor $q+1$.
Different encodings may compute the same function, which only helps the
upper bound.
\end{proof}

\begin{theorem}[Hard powers of a simple permutation]\label{thm:lower}
For all sufficiently large even $n$, the permutation in
\cref{eq:permutation} satisfies
\begin{equation}\label{eq:hard-power}
  \max_{T\geq 0}\CC(P_n^T)
    =\Omega\!\left(\frac{2^{n/2}}{n}\right).
\end{equation}
Together with \cref{thm:main}, this gives
\begin{equation}\label{eq:tight-exponent}
  \log_2\!\left(\max_{T\geq 0}\CC(P_n^T)\right)
    =\frac n2+O(\log n).
\end{equation}
\end{theorem}

\begin{proof}
There are exactly $D_B$ distinct powers of $P_n$, with representatives
$0\leq T<D_B$.  By \cref{lem:lcm},
\[
  \log_2 D_B\geq B-\log_2(B+1).
\]
Let $q=\floor{cB/n}$ for a sufficiently small positive absolute constant
$c$.  Since $B=2^{n/2}-1$, \cref{lem:count} bounds the logarithm of the
number of circuits of size at most $q$ by
\[
  O((q+n)\log(n+q+2))=O(cB+n^2).
\]
Choose $c$ so that the coefficient of $B$ is less than $1/2$.
For all sufficiently large even $n$, this quantity is smaller than
$B-\log_2(B+1)$.  Thus there are fewer size-$q$ circuits than distinct
powers of $P_n$.  At least one power has complexity greater than $q$,
proving \cref{eq:hard-power}.

For the upper bound, substitute $\CC(P_n)=O(n)$ into
\cref{eq:optimized}.  This gives
$\max_T\CC(P_n^T)=O(n^3 2^{n/2})$.  Combining the upper and lower
bounds and taking logarithms proves \cref{eq:tight-exponent}.
\end{proof}

\begin{corollary}\label[corollary]{cor:one-output}
For all sufficiently large even $n$, some power of $P_n$ has an output
coordinate of circuit complexity $\Omega(2^{n/2}/n^2)$.
\end{corollary}

\begin{proof}
For the power supplied by \cref{thm:lower}, join optimal circuits for all
$n$ output coordinates.  Their total size is at most $n$ times the
largest coordinate complexity and is an upper bound on the complexity
of the whole map.
\end{proof}

\subsection{Why the lower bound does not apply to short input times}
\label{sec:lower-limits}

The lower bound selects $T$ from $\{0,\ldots,D_B-1\}$; it places no
polynomial bound on the bit length of that exponent.  Indeed,
\cref{lem:lcm} shows that $\log D_B\geq 2^{n/2}-O(n)$.
The choice of a power therefore has exponentially many available
bits.  The theorem does not provide a polynomial-time procedure that,
given $n$, identifies a hard power.

For the very same family, the iterator with an $m$-bit time input has
size $O(nm+n^2)$.  On an active cycle,
\begin{equation}\label{eq:easy-iterator}
  P_n^t(a,b)=(a,(b+(t\bmod a))\bmod a).
\end{equation}
An $m$-stage remainder computation with a $k$-bit running remainder costs
$O(mk)$ gates, and the remaining comparison, addition, and selection
cost at most $O(n^2)$ gates.  On the inactive states the answer is the
input itself; use a dummy positive modulus when $a=0$.
Consequently, for $m=\poly(n)$ this iterator has polynomial size.

\begin{corollary}[Bit length of hard exponents]\label[corollary]{cor:exponent-length}
Let $T_n$ be exponents satisfying
$\CC(P_n^{T_n})=\Omega(2^{n/2}/n)$ for sufficiently large even $n$.
Their binary lengths
$m_n=\max\{1,\ceil{\log_2(T_n+1)}\}$ satisfy
\[
  m_n=\Omega(2^{n/2}/n^2).
\]
\end{corollary}

\begin{proof}
Hardwire the $m_n$-bit exponent into the iterator just described.
This gives $\CC(P_n^{T_n})=O(nm_n+n^2)$.
Comparing with the assumed lower bound and using
$n^2=o(2^{n/2}/n)$ proves the claim.
\end{proof}

Thus the hard powers necessarily require exponentially long binary
exponents; long descriptions are not merely allowed by the counting proof.
Theorem~\ref{thm:lower} establishes sharpness for arbitrary fixed powers.
It gives no hardness result for the bounded-bit-time problem.

\section{Permutation inversion and one-way functions}\label{sec:one-way}

\subsection{Inversion as a fixed power}

\begin{corollary}[Exact permutation inversion]\label[corollary]{cor:inverse}
Let $P:\bits^n\to\bits^n$ be a permutation with $\CC(P)\leq s$, and put
$N=2^n$ and $A=s+n^3$.  For every $1\leq L\leq N$,
\[
  \CC(P^{-1})
  =O\!\left(Ls+n^3\left(L+\frac NL\right)\right).
\]
In particular,
\[
  \CC(P^{-1})=O\!\left(A+\sqrt{n^3NA}\right).
\]
If $s=n^{O(1)}$, an inverter of size $2^{n/2+O(\log n)}$ is correct on
every input.
\end{corollary}

\begin{proof}
Let $D=\ord(P)$, the least common multiple of its cycle lengths.
Then $P^D$ is the identity, so $P^{-1}=P^{D-1}$.
Apply \cref{thm:main}.  The exponent may have a long binary expansion,
but the theorem charges only the resulting tables and their access gates.
\end{proof}

This recovers the permutation-inversion circuit bound already discussed
by Hirahara, Ilango, and Williams~\cite[Section~3.2]{HIW24}.
It identifies an earlier special case of the fixed-power theorem.

\subsection{Compatibility with one-wayness}

A cryptographic one-way function is computable in polynomial time, while
every probabilistic polynomial-time inverter succeeds with only negligible
probability on the image of a uniform input.  One may instead require
security against every polynomial-size circuit family, a stronger
nonuniform notion~\cite[Section~7.1.1]{Goldreich08}.  For a length-preserving
family $f_n:\bits^n\to\bits^n$, the latter requires that, for every such
family $C_n$,
\begin{equation}\label{eq:one-way}
  \Pr_{x\leftarrow\bits^n}
  \bigl[f_n(C_n(f_n(x)))=f_n(x)\bigr]
\end{equation}
be negligible in $n$: smaller than $n^{-c}$ eventually for every fixed
$c>0$.  Any preimage counts as success.  Deterministic circuits suffice
in this definition, since random coins can be fixed to maximize average
success.  A one-way permutation is the bijective special case.

The guarantee in \cref{cor:inverse} is an exponential size upper bound.
It therefore does not contradict either polynomial-time or polynomial-size
one-wayness.  Nonuniformity alone is not the reason for compatibility:
\cref{eq:one-way} already allows arbitrary nonuniform advice encoded in
polynomial-size circuits.  The relevant distinction is the circuit size.
Moreover, no efficient procedure for constructing our inverter is asserted.

There is a quantitative consequence.  For every fixed $\varepsilon>0$,
every polynomial-size permutation family on $n$ bits has exact inverters
of size at most $2^{(1/2+\varepsilon)n}$ for all sufficiently large $n$.
Thus such a family cannot remain one-way against all circuits of that
size, with any success bound strictly below one.  This upper bound leaves
open security against polynomial-size circuits, and does not exclude
security against size $2^{cn}$ for a fixed $0<c<1/2$.  It does not
establish that either kind of security is achievable.

\subsection{General functions and the counting example}

For a noninjective map, inversion is not generally a forward power.
For example, let $f$ on $\{0,1,2,3\}$ have values $(0,0,1,0)$.
The unique preimage of $1$ is $2$, whereas every forward iterate of $1$
belongs to $\{0,1\}$.  Thus \cref{thm:main} does not directly provide
square-root-size inverters for arbitrary one-way functions.
For comparison, Hirahara, Ilango, and Williams prove a general inversion
bound of $2^{4n/5}\poly(s+n)$ gates~\cite[Theorem~3.5]{HIW24}.

Encoding exhaustive preimage search as iteration does not avoid the state
accounting.  The straightforward transition retains an $n$-bit target
$y$, an $n$-bit candidate $x$, and a flag; it tests $f(x)=y$ and either
halts or increments $x$.  This uses $2n+O(1)$ state bits.  For
polynomial-size $f$, our theorem then gives size $2^{n+O(\log n)}$,
not the half exponent measured in the preimage length.

The counting construction in \cref{sec:lower} also does not establish
one-way functions.  Its update permutation $P_n$ has an $O(n)$-gate
inverse: on an active cycle decrement $b$ modulo $a$, and leave inactive
states fixed.  A family of the selected hard powers has exponential
worst-case \emph{forward} circuit complexity, so it cannot be a
polynomial-time one-way function.  The counting proof supplies no
average-case inversion hardness.  The results here consequently neither
prove nor refute the existence of one-way functions or one-way permutations.

\section{Scope and remaining questions}\label{sec:scope}

The main theorem is an exact nonuniform upper bound for unrestricted
bounded-fan-in circuits.  It applies to noninvertible maps and is uniform
in the starting state, but not in the circuit construction or its
preprocessing cost.  It gives no depth bound and no efficient procedure
for discovering checkpoints or their payloads from a succinct
representation of $F$.  An initial-state-dependent family of update
rules must be encoded as a single common transition map before the
theorem applies; any additional persistent state changes $n$ and hence
the exponential scale of the bound.

The sharpness result deliberately separates the simplicity of an update
rule from the description complexity of a chosen power.  A stronger
lower bound for a specified iteration count, or for an iterator with an
$n$-bit time input, is not established here.  The explicit permutation
used for counting cannot witness such a strengthening, by
\cref{eq:easy-iterator,cor:exponent-length}.  The inversion consequence
in \cref{cor:inverse} likewise leaves polynomial-size one-wayness open.

Two questions remain distinct.  First, can stronger upper bounds hold
when the exponent is restricted to polynomially many bits or when the
functional graph has additional structure?  Second, under what conditions
can the required preprocessing be performed substantially faster than
exploring the full state space?  The present proof addresses neither
question.  Its conclusion is the unconditional total-gate tradeoff:
a short transition circuit can be combined with a sparse table, whose
access cost is explicitly accounted for, to compute every fixed iterate
at square-root state-space scale.

\section*{Provenance and draft status}

This draft and the underlying argument were developed with assistance
from ChatGPT (GPT-6 Astra Pro), including assistance with theorem
formulation, proof development, exposition, and literature searches.
The author field is provisional.  Historical priority and final
attribution require author review; the draft does not represent
independent peer verification.

\begin{thebibliography}{9}

\bibitem{Batcher68}
Kenneth E. Batcher.
\newblock Sorting networks and their applications.
\newblock In \emph{Proceedings of the April 30--May 2, 1968, Spring Joint
Computer Conference}, AFIPS Conference Proceedings, volume 32,
pages 307--314, 1968.
\newblock \href{https://doi.org/10.1145/1468075.1468121}{doi:10.1145/1468075.1468121}.

\bibitem{Hellman80}
Martin E. Hellman.
\newblock A cryptanalytic time--memory trade-off.
\newblock \emph{IEEE Transactions on Information Theory}, 26(4):401--406, 1980.
\newblock \href{https://ee.stanford.edu/~hellman/publications/36.pdf}{Author-hosted full text}.

\bibitem{MRRR12}
J. Ian Munro, Rajeev Raman, Venkatesh Raman, and S. Srinivasa Rao.
\newblock Succinct representations of permutations and functions.
\newblock \emph{Theoretical Computer Science}, 438:74--88, 2012.
\newblock Preprint: \href{https://arxiv.org/abs/1108.1983}{arXiv:1108.1983}, 2011.

\bibitem{HIW24}
Shuichi Hirahara, Rahul Ilango, and R. Ryan Williams.
\newblock Beating brute force for compression problems.
\newblock In \emph{Proceedings of the 56th Annual ACM Symposium on Theory
of Computing (STOC)}, pages 659--670, 2024.
\newblock \href{https://doi.org/10.1145/3618260.3649778}{doi:10.1145/3618260.3649778}.
\newblock Full version: \href{https://eccc.weizmann.ac.il/report/2023/171/}{ECCC TR23-171};
section and theorem numbers refer to that version.

\bibitem{Eppstein23}
David Eppstein.
\newblock The complexity of iterated reversible computation.
\newblock \emph{TheoretiCS}, 2, Article 10, 1--41, 2023.
\newblock \href{https://doi.org/10.46298/theoretics.23.10}{doi:10.46298/theoretics.23.10}.

\bibitem{Goldreich08}
Oded Goldreich.
\newblock \emph{Computational Complexity: A Conceptual Perspective}.
\newblock Cambridge University Press, 2008.
\newblock Author's draft, Chapter~7:
\href{https://www.wisdom.weizmann.ac.il/~oded/CC/r3.pdf}{one-way function definitions}.

\end{thebibliography}

\end{document}
